\documentclass[a4paper]{article}
\input{preamble.tex}
\title{Analysis 2: HW6}
\author{Aamod Varma}
\begin{document}
\maketitle
% \textbf{Problem 1.}
% We have $\Delta u = \sum_{i= 1}^{n} \partial_{ii} u = 0$

% \qquad 1. First we show that $u(x) = \frac{1}{|\partial B(x, r) |} \int_{\partial B(x, r)}^{}  u \ dS$. So consider,
% \[
%     \frac{1}{|\partial B(x, r) |} \int_{\partial B(x, r)}^{}  u \ dS
% \]

% and we can write this as, 


% \[
%     \int_{\partial B(x, \rho)}^{} u \ dS = \int_{\partial B(0, 1)}^{} u (x + r \theta) r^{n - 1} dS(\theta)
% \]

% Now differentiating and applying $\Delta u = 0$ we can get, 
% \[
%     u(x) = \frac{1}{\left |\partial B(x, r) \right |} \int_{\partial B(x, r)}^{} u \ dS
% \]

% now we need to show the volume average part, so we can write the volume integral as, 
% \begin{align*}
%     \int_{B(x,r)}^{}  u = \int_{0}^{r} \left ( \int_{\partial B(x, \rho)}^{} u \ dS\right ) \ d \rho
% \end{align*}

% but note that we have, 
% \[
%     \int_{\partial B(x, \rho)}^{} u dS = \left | \partial B(x, \rho)\right | u(x)
% \]

% putting that in we get, 
% \[
%     \int_{B(x, r)}^{} u = u(x) \int_{0}^{r} \left | \partial B(x, \rho)\right | d \rho = u(x) \left | B(x, r)\right |
% \]

% which gives us, 
% \[
%     u(x) = \frac{1}{\left | B(x, r)\right |} \int_{B(x, r)}^{} u
% \]


% \vspace{1em}

% 2.  We have $\phi$ such that $\int_{R^{n}}^{}  \phi = 1$. So,
% \begin{align*}
%     u(x) = (u * \phi)(x) = \int_{R^{n}}^{} u(x - y)\phi(y) \ dy
% \end{align*}

% Now given $\Delta (u * \phi)(x)$ and  that $\Delta u = 0$ we can apply $\Delta$ to both to get $\Delta (u * \phi) = 0$. Now apply (1) to the function $f = (u * \phi)$ to get,
% \[
%     f(x) = \int_{}^{} u(x - y) \phi(y) dy = u(x)
% \]

% so this means $f = u * \phi = u$. And since $u * \phi$ is smooth, we have $u$ is smooth.


% \vspace{1em}

% \textbf{Problem 2}
% \qquad  Choose $x \ne y$ and take $u(x), u(y)$ as, 
% \[
%     u(x) = \frac{1}{\left | \partial B (x, r)\right |} \int_{\partial B(x, r)}^{} u(x) d S
% \]
% and, 

% \[
%     u(y) = \frac{1}{\left | \partial B (y, r)\right |} \int_{\partial B(y, r)}^{} u(y) d S
% \]

% now choose large enough $r$ relative to $x$ and $y$ so that the surfaces of the balls centered at $x$ and $y$ are basically the same. Now note that, 
% \[
%     |u(x) - u(y)| = \left | \frac{1}{\left | \partial B (x, r)\right |} \int_{\partial B(x, r)}^{} u(x)  -  \frac{1}{\left | \partial B (y, r)\right |} \int_{\partial B(y, r)}^{} u(y)  \right |
% \]

% now note that this can be bounded inversely by $r$  so we have, 
% \[
%     |u(x)  - u(y) | \le \frac{C}{r}
% \]

% for some constant $C$. So if $r$ is very large we have this goes to zero and $u(x) = u(y)$


% \vspace{1em}

\textbf{Problem 3}.
We have $u$ is continuous and that, 
\[
    \int_{R^{n}}^{}  u \Delta \psi = 0
\]

we need to show that $u$ is smooth and $\Delta u = 0$.
\vspace{1em}

We can approximate $u$ using the convolution between $u$ and $\phi$ where we define $\phi = \frac{1}{\epsilon^{n}} f$ where $\int_{}^{}f = 1 $. And now note that we  get $u_{\epsilon} = u * \phi$ but if we take $n$ to infinity then we have the approximation of $u_{\epsilon} \to u$. Now note that, 
\[
    \Delta u_{\epsilon} = u * \Delta \phi = \int_{R^{n}}^{}  u(y) \Delta \phi(x - y) dy
\]

but the RHS is zero so we have $\Delta u_{\epsilon} = 0$. But now use the notating from (1) to get, 
\[
    u_{\epsilon}(x) = \frac{1}{\left | B(x, r)\right |} \int_{B(x, r)}^{}  u_{\epsilon}(y) dy
\]

and take $\epsilon \to 0$ then as we have $u_{\epsilon} \to 0$ we have $u$ also gets the same property.

\vspace{1em}

So we showed that $u$ is smooth and  $\Delta u = 0$

\end{document}

